% questions.tex —— 2014 年·湖南卷·理科 · 题目内容片段（无答案）
% 由 tools/generate_projects.py 依来源总账自 2014/2014_数学真题_湖南理_试题.md 生成，勿手改（改源 MD 后重跑）


\begin{gknotice}
\textbf{来源说明}：本题面依据固定源 PDF 的两页 A3 题面逐题转录，原卷印刷页码为 885—886。第 885 面为选择题 1—10、填空题 11—16 与解答题 17—18；第 886 面为解答题 19—22。全卷三处节头依次为「一、选择题」「二、填空题」「三、解答题」。

\textbf{解析归属}：固定源不附答案或解析；配套答案与解析均由本库按本题面独立复算。

\textbf{判读备注}：第 1 题分子印面上 $z$ 的右上角有一处短横状墨痕；按 $\bar z$ 解读时方程 $\dfrac{\bar z+\mathrm i}{z}=\mathrm i$ 无解（比较实虚部得 $1=0$），按 $z+\mathrm i$ 解读则唯一对应选项（B），故题面取 $z+\mathrm i$。

\textbf{图形说明}：第 6 题的程序框图、第 7 题的三视图、第 12 题的圆与弦、第 15 题的正方形与抛物线、第 18 题的平面四边形、第 19 题的四棱柱、第 21 题的椭圆与双曲线由项目内 TikZ 矢量源重绘；其余各题原页无图，不补绘。
\end{gknotice}

\gkpart{一、选择题}

\begin{question}
\gkti{1}满足 $\dfrac{z+\mathrm i}{z}=\mathrm i$（$\mathrm i$ 为虚数单位）的复数 $z=$\mbox{（　　）}

（A）$\dfrac12+\dfrac12\mathrm i$；（B）$\dfrac12-\dfrac12\mathrm i$；（C）$-\dfrac12+\dfrac12\mathrm i$；（D）$-\dfrac12-\dfrac12\mathrm i$。

\end{question}
\begin{question}
\gkti{2}对一个容量为 $N$ 的总体抽取容量为 $n$ 的样本，当选取简单随机抽样、系统抽样和分层抽样三种不同方法抽取样本时，总体中每个个体被抽中的概率分别为 $p_1$，$p_2$，$p_3$，则\mbox{（　　）}

（A）$p_1=p_2<p_3$；（B）$p_2=p_3<p_1$；（C）$p_1=p_3<p_2$；（D）$p_1=p_2=p_3$。

\end{question}
\begin{question}
\gkti{3}已知 $f(x)$，$g(x)$ 分别是定义在 $\mathbf R$ 上的偶函数和奇函数，且 $f(x)-g(x)=x^3+x^2+1$，则 $f(1)+g(1)=$\mbox{（　　）}

（A）$-3$；（B）$-1$；（C）$1$；（D）$3$。

\end{question}
\begin{question}
\gkti{4}$\left(\dfrac12x-2y\right)^5$ 的展开式中 $x^2y^3$ 的系数是\mbox{（　　）}

（A）$-20$；（B）$-5$；（C）$5$；（D）$20$。

\end{question}
\begin{question}
\gkti{5}已知命题 $p$：若 $x>y$，则 $-x<-y$；命题 $q$：若 $x>y$，则 $x^2>y^2$。在命题 ① $p\wedge q$；② $p\vee q$；③ $p\wedge(\neg q)$；④ $(\neg p)\vee q$ 中，真命题是\mbox{（　　）}

（A）①③；（B）①④；（C）②③；（D）②④。

\Needspace{0.4\textheight}
\end{question}
\begin{question}
\gkti{6}执行如图所示的程序框图，如果输入的 $t\in[-2,2]$，则输出的 $S$ 属于\mbox{（　　）}
\nopagebreak[4]
\begin{center}% ============================================================
% f2014_hunan_science_q06_flowchart.tex
% 2014 湖南卷（理科）第 6 题：程序框图
% 来源 PDF：普通高考/2014/2014湖南理.pdf
% 原页：第 1 页（A3 横排）第 1 栏下部，印面 885，第 6 题题干下方
% 结构：开始 → 输入 t → 判断 t<0；
%       是：t=2t^2+1，回到判断框；否：S=t-3 → 输出 S → 结束。
% ============================================================
\begin{tikzpicture}[>=Stealth, line width=0.7pt, font=\small,
  x=1cm, y=1cm, line cap=round, line join=round]
  % 开始
  \draw[rounded corners=6pt] (-0.80,0.35) rectangle (0.80,-0.35);
  \node at (0,0) {开始};
  \draw[->] (0,-0.35) -- (0,-0.62);
  % 输入 t
  \draw (-1.15,-0.62) -- (1.15,-0.62) -- (1.00,-1.28) -- (-1.30,-1.28) -- cycle;
  \node at (0,-0.95) {输入 $t$};
  \draw[->] (0,-1.28) -- (0,-1.58);
  % 判断框
  \draw (0,-1.58) -- (1.95,-2.28) -- (0,-2.98) -- (-1.95,-2.28) -- cycle;
  \node at (0,-2.28) {$t<0$?};
  % 是：向左再向下进入循环体
  \draw[->] (-1.95,-2.28) -- (-2.85,-2.28) -- (-2.85,-3.25);
  \node at (-2.35,-2.12) {是};
  \draw (-3.80,-3.25) rectangle (-1.90,-3.85);
  \node at (-2.85,-3.55) {$t=2t^2+1$};
  % 回边：循环体左侧向上，水平进入判断框左上边
  \draw[->] (-3.80,-3.55) -- (-4.45,-3.55) -- (-4.45,-1.80) -- (-0.61,-1.80);
  % 否：向右再向下
  \draw[->] (1.95,-2.28) -- (3.00,-2.28) -- (3.00,-3.12);
  \node at (2.42,-2.12) {否};
  \draw (2.05,-3.12) rectangle (3.95,-3.72);
  \node at (3.00,-3.42) {$S=t-3$};
  % 输出 S
  \draw[->] (3.00,-3.72) -- (3.00,-4.25) -- (0,-4.25) -- (0,-4.45);
  \draw (-1.35,-4.45) -- (1.35,-4.45) -- (1.05,-5.15) -- (-1.65,-5.15) -- cycle;
  \node at (-0.30,-4.80) {输出 $S$};
  \draw[->] (0,-5.15) -- (0,-5.45);
  % 结束
  \draw[rounded corners=6pt] (-0.80,-5.45) rectangle (0.80,-6.15);
  \node at (0,-5.80) {结束};
\end{tikzpicture}
\end{center}
\nopagebreak[4]
（A）$[-6,-2]$；（B）$[-5,-1]$；（C）$[-4,5]$；（D）$[-3,6]$。

\Needspace{0.4\textheight}
\end{question}
\begin{question}
\gkti{7}一块石材表示的几何体的三视图如图所示，将该石材切削、打磨，加工成球，则能得到的最大球的半径等于\mbox{（　　）}
\nopagebreak[4]
\begin{center}% ============================================================
% f2014_hunan_science_q07_views.tex
% 2014 湖南卷（理科）第 7 题：石材几何体的三视图
% 来源 PDF：普通高考/2014/2014湖南理.pdf
% 原页：第 1 页（A3 横排）第 1 栏底部，印面 885，第 7 题题干下方
% 结构：正视图为直角边长 6、8 的直角三角形；侧视图为矩形；
%       俯视图为 6×12 的矩形，标注柱长 12。比例 0.4 cm/单位。
% ============================================================
\begin{tikzpicture}[>=Stealth, line width=0.6pt, font=\small, scale=0.72,
  x=1cm, y=1cm, line cap=round, line join=round]
  % 正视图：直角边 6（横）、8（纵）
  \draw (0,0) -- (2.4,0) -- (0,3.2) -- cycle;
  \draw[<->] (0,3.95) -- (2.4,3.95);
  \node at (1.2,4.22) {$6$};
  \draw[<->] (-0.95,0) -- (-0.95,3.2);
  \node at (-1.18,1.6) {$8$};
  \node at (1.2,-0.52) {正视图};
  % 侧视图：12×8
  \draw (4.4,0) rectangle (9.2,3.2);
  \node at (6.8,-0.52) {侧视图};
  % 俯视图：6×12
  \draw (10.4,-1.6) rectangle (12.8,3.2);
  \draw[<->] (13.10,-1.6) -- (13.10,3.2);
  \node at (13.52,0.8) {$12$};
  \node at (11.6,-2.12) {俯视图};
\end{tikzpicture}
\end{center}
\nopagebreak[4]
（A）$1$；（B）$2$；（C）$3$；（D）$4$。

\end{question}
\begin{question}
\gkti{8}某市生产总值连续两年持续增加。第一年的增长率为 $p$，第二年的增长率为 $q$，则该市这两年生产总值的年平均增长率为\mbox{（　　）}

（A）$\dfrac{p+q}2$；（B）$\dfrac{(p+1)(q+1)-1}2$；（C）$\sqrt{pq}$；（D）$\sqrt{(p+1)(q+1)}-1$。

\end{question}
\begin{question}
\gkti{9}已知函数 $f(x)=\sin(x-\varphi)$，且 $\displaystyle\int_0^{\frac{2\pi}3}f(x)\,\mathrm dx=0$，则函数 $f(x)$ 的图象的一条对称轴是\mbox{（　　）}

（A）$x=\dfrac{5\pi}6$；（B）$x=\dfrac{7\pi}{12}$；（C）$x=\dfrac\pi3$；（D）$x=\dfrac\pi6$。

\end{question}
\begin{question}
\gkti{10}已知函数 $f(x)=x^2+\mathrm e^x-\dfrac12$（$x<0$）与 $g(x)=x^2+\ln(x+a)$ 的图象上存在关于 $y$ 轴对称的点，则 $a$ 的取值范围是\mbox{（　　）}

（A）$\left(-\infty,\dfrac1{\sqrt{\mathrm e}}\right)$；（B）$(-\infty,\sqrt{\mathrm e})$；（C）$\left(-\dfrac1{\sqrt{\mathrm e}},\sqrt{\mathrm e}\right)$；（D）$\left(-\sqrt{\mathrm e},\dfrac1{\sqrt{\mathrm e}}\right)$。

\end{question}
\gkpart{二、填空题}

\begin{question}
\gkti{11}在平面直角坐标系中，倾斜角为 $\dfrac\pi4$ 的直线 $l$ 与曲线 $C:\begin{cases}x=2+\cos\alpha,\\y=1+\sin\alpha,\end{cases}$（$\alpha$ 为参数）交于 $A$，$B$ 两点，且 $|AB|=2$，以坐标原点 $O$ 为极点、$x$ 轴正半轴为极轴建立极坐标系，则直线 $l$ 的极坐标方程是\gkblank{}。

\Needspace{0.4\textheight}
\end{question}
\begin{question}
\gkti{12}如图，已知 $AB$，$BC$ 是 $\odot O$ 的两条弦，$AO\perp BC$，$AB=\sqrt3$，$BC=2\sqrt2$，则 $\odot O$ 的半径等于\gkblank{}。
\nopagebreak[4]
\begin{center}% ============================================================
% f2014_hunan_science_q12_circle.tex
% 2014 湖南卷（理科）第 12 题：圆 O 的两条弦 AB、BC
% 来源 PDF：普通高考/2014/2014湖南理.pdf
% 原页：第 1 页（A3 横排）第 2 栏中部，印面 885，第 12 题题干下方
% 结构：O 为圆心，半径 R=3/2；A(-3/2,0) 在圆上，AO 为半径；
%       BC 为垂直于 AO 的弦，B、C 分别在 AO 上下，BM=MC=√2，
%       OM=1/2。比例 1.3 cm/单位。
% ============================================================
\begin{tikzpicture}[>=Stealth, line width=0.7pt, font=\small,
  x=1.3cm, y=1.3cm, line cap=round, line join=round]
  % 圆
  \draw (0,0) circle (1.5);
  % 半径 AO
  \draw (-1.5,0) -- (0,0);
  % 弦 BC（过 AO 上的点 M）
  \draw (-0.5,1.41421) -- (-0.5,-1.41421);
  % 弦 AB
  \draw (-1.5,0) -- (-0.5,1.41421);
  % 圆心与顶点
  \fill (0,0) circle (1.4pt);
  \node[right] at (0.06,0.02) {$O$};
  \node[left]  at (-1.56,0.10) {$A$};
  \node[above] at (-0.52,1.52) {$B$};
  \node[below] at (-0.52,-1.52) {$C$};
\end{tikzpicture}
\end{center}
\nopagebreak[4]
\end{question}
\begin{question}
\gkti{13}若关于 $x$ 的不等式 $|ax-2|<3$ 的解集为 $\left\{x\mid-\dfrac53<x<\dfrac13\right\}$，则 $a=$\gkblank{}。

\end{question}
\begin{question}
\gkti{14}若变量 $x$，$y$ 满足约束条件 $\begin{cases}y\leqslant x,\\x+y\leqslant4,\\y\geqslant k,\end{cases}$ 且 $z=2x+y$ 的最小值为 $-6$，则 $k=$\gkblank{}。

\Needspace{0.4\textheight}
\end{question}
\begin{question}
\gkti{15}如图，正方形 $ABCD$ 和正方形 $DEFG$ 的边长分别为 $a$，$b$（$a<b$），原点 $O$ 为 $AD$ 的中点，抛物线 $y^2=2px$（$p>0$）经过 $C$，$F$ 两点，则 $\dfrac ba=$\gkblank{}。
\nopagebreak[4]
\begin{center}% ============================================================
% f2014_hunan_science_q15_squares_parabola.tex
% 2014 湖南卷（理科）第 15 题：正方形 ABCD、DEFG 与抛物线 y^2=2px
% 来源 PDF：普通高考/2014/2014湖南理.pdf
% 原页：第 1 页（A3 横排）第 3 栏顶部，印面 885，第 15 题题干上方
% 结构：O 为原点，AD 在 x 轴上且 O 为 AD 中点，|AD|=a=1；
%       ABCD 在 x 轴下方，DEFG 在 x 轴上方且以 D 为左顶点，|DE|=b=1+√2；
%       C、F 在抛物线 y^2=2px（p=a=1）上。比例 1.2 cm/单位。
% ============================================================
\begin{tikzpicture}[>=Stealth, line width=0.6pt, font=\small,
  x=1.2cm, y=1.2cm, line cap=round, line join=round]
  % 坐标轴
  \draw[->] (-1.05,0) -- (3.75,0) node[right] {$x$};
  \draw[->] (0,-2.45) -- (0,2.90) node[above] {$y$};
  % 抛物线 y^2=2x
  \draw[domain=-2.65:2.65, samples=150, smooth, variable=\t]
    plot ({(\t*\t)/2}, {\t});
  % 正方形 ABCD（x 轴下方）
  \draw (-0.5,0) -- (-0.5,-1) -- (0.5,-1) -- (0.5,0) -- cycle;
  % 正方形 DEFG（x 轴上方）
  \draw (0.5,0) -- (2.91421,0) -- (2.91421,2.41421) -- (0.5,2.41421) -- cycle;
  % 原点与顶点标注
  \node[below left]  at (-0.06,-0.04) {$O$};
  \node[below left]  at (-0.5,-0.06) {$A$};
  \node[below left]  at (-0.5,-1.02) {$B$};
  \node[right]       at (0.56,-0.96) {$C$};
  \node[below right] at (0.5,-0.06) {$D$};
  \node[below]       at (2.91421,-0.10) {$E$};
  \node[above right] at (2.88000,2.47000) {$F$};
  \node[above left]  at (0.52000,2.47000) {$G$};
\end{tikzpicture}
\end{center}
\nopagebreak[4]
\end{question}
\begin{question}
\gkti{16}在平面直角坐标系中，$O$ 为原点，$A(-1,0)$，$B(0,\sqrt3)$，$C(3,0)$，动点 $D$ 满足 $|\overrightarrow{CD}|=1$，则 $|\overrightarrow{OA}+\overrightarrow{OB}+\overrightarrow{OD}|$ 的最大值是\gkblank{}。

\end{question}
\gkpart{三、解答题}

\begin{question}
\gkti{17}某企业有甲、乙两个研发小组，他们研发新产品成功的概率分别为 $\dfrac23$ 和 $\dfrac35$。现安排甲组研发新产品 $A$，乙组研发新产品 $B$。设甲、乙两组的研发相互独立。

（1） 求至少有一种新产品研发成功的概率；

（2） 若新产品 $A$ 研发成功，预计企业可获利润 $120$ 万元；若新产品 $B$ 研发成功，预计企业可获利润 $100$ 万元。求该企业可获利润的分布列和数学期望。

\Needspace{0.4\textheight}
\end{question}
\begin{question}
\gkti{18}如图，在平面四边形 $ABCD$ 中，$AD=1$，$CD=2$，$AC=\sqrt7$。
\nopagebreak[4]
\begin{center}% ============================================================
% f2014_hunan_science_q18_quad.tex
% 2014 湖南卷（理科）第 18 题：平面四边形 ABCD 与对角线 AC
% 来源 PDF：普通高考/2014/2014湖南理.pdf
% 原页：第 1 页（A3 横排）第 3 栏下部，印面 885，第 18 题题干下方
% 结构：按题设精确作图（比例 0.75 cm/单位）：
%       A(0,0)，D(1,0)，C(2,-√3)，B(-0.616,-3.200)；
%       即 AD=1，CD=2，AC=√7，∠BAD≈100.9°，BC=3。
% ============================================================
\begin{tikzpicture}[>=Stealth, line width=0.7pt, font=\small,
  x=1cm, y=1cm, line cap=round, line join=round]
  \coordinate (A) at (1.00,2.55);
  \coordinate (B) at (0.538,0.150);
  \coordinate (C) at (2.500,1.251);
  \coordinate (D) at (1.750,2.550);
  % 四边形边界
  \draw (A) -- (B);
  \draw (B) -- (C);
  \draw (C) -- (D);
  \draw (D) -- (A);
  % 对角线 AC
  \draw (A) -- (C);
  % 顶点标注
  \node[above]       at (A) {$A$};
  \node[below left]  at (B) {$B$};
  \node[below right] at (C) {$C$};
  \node[above right] at (D) {$D$};
\end{tikzpicture}
\end{center}
\nopagebreak[4]
（1） 求 $\cos\angle CAD$ 的值；

（2） 若 $\cos\angle BAD=-\dfrac{\sqrt7}{14}$，$\sin\angle CBA=\dfrac{\sqrt{21}}6$，求 $BC$ 的长。

\Needspace{0.4\textheight}
\end{question}
\begin{question}
\gkti{19}如图，四棱柱 $ABCD-A_1B_1C_1D_1$ 的所有棱长都相等，$AC\cap BD=O$，$A_1C_1\cap B_1D_1=O_1$，四边形 $ACC_1A_1$ 和四边形 $BDD_1B_1$ 均为矩形。
\nopagebreak[4]
\begin{center}% ============================================================
% f2014_hunan_science_q19_prism.tex
% 2014 湖南卷（理科）第 19 题：四棱柱 ABCD-A1B1C1D1
% 来源 PDF：普通高考/2014/2014湖南理.pdf
% 原页：第 2 页（A3 横排）左栏，印面 886，第 19 题题干下方
% 结构：底面为菱形 ABCD（A 在后、B 在左、C 在前、D 在右），
%       O、O1 分别为 AC∩BD 与 A1C1∩B1D1；侧棱垂直于底面。
%       实线：顶面四边、三条可见侧棱、BC、CD；
%       虚线：AA1、AB、AD、AC、BD、OO1、OB1、OC1。
% ============================================================
\begin{tikzpicture}[>=Stealth, line width=0.7pt, font=\small,
  x=1cm, y=1cm, line cap=round, line join=round]
  \coordinate (A)  at (1.30,1.70);
  \coordinate (B)  at (0.00,0.55);
  \coordinate (C)  at (2.10,0.00);
  \coordinate (D)  at (3.40,1.15);
  \coordinate (A1) at (1.30,4.10);
  \coordinate (B1) at (0.00,2.95);
  \coordinate (C1) at (2.10,2.40);
  \coordinate (D1) at (3.40,3.55);
  \coordinate (O)  at (1.70,0.85);
  \coordinate (O1) at (1.70,3.25);
  % 顶面
  \draw (A1) -- (B1);
  \draw (B1) -- (C1);
  \draw (C1) -- (D1);
  \draw (D1) -- (A1);
  % 可见侧棱
  \draw (B1) -- (B);
  \draw (C1) -- (C);
  \draw (D1) -- (D);
  % 底面可见边
  \draw (B) -- (C);
  \draw (C) -- (D);
  % 底面与内部虚线
  \draw[dashed] (A1) -- (A);
  \draw[dashed] (A) -- (B);
  \draw[dashed] (A) -- (D);
  \draw[dashed] (A) -- (C);
  \draw[dashed] (B) -- (D);
  \draw[dashed] (O) -- (O1);
  \draw[dashed] (O) -- (B1);
  \draw[dashed] (O) -- (C1);
  % 顶点标注
  \node[above left,  inner sep=2pt] at (A1) {$A_1$};
  \node[left,        inner sep=2pt] at (B1) {$B_1$};
  \node[right,       inner sep=2pt] at (C1) {$C_1$};
  \node[above right] at (D1) {$D_1$};
  \node[inner sep=2pt] at (1.10,1.30) {$A$};
  \node[below left,  inner sep=2pt] at (B)  {$B$};
  \node[below,       inner sep=2pt] at (C)  {$C$};
  \node[right,       inner sep=2pt] at (D)  {$D$};
  \node[below,       inner sep=2pt] at (O)  {$O$};
  \node[above,       inner sep=2pt] at (O1) {$O_1$};
\end{tikzpicture}
\end{center}
\nopagebreak[4]
（1） 证明：$O_1O\perp$ 底面 $ABCD$；

（2） 若 $\angle CBA=60^\circ$，求二面角 $C_1-OB_1-D$ 的余弦值。

\end{question}
\begin{question}
\gkti{20}已知数列 $\{a_n\}$ 满足 $a_1=1$，$|a_{n+1}-a_n|=p^n$，$n\in\mathbf N^*$。

（1） 若 $\{a_n\}$ 是递增数列，且 $a_1$，$2a_2$，$3a_3$ 成等差数列，求 $p$ 的值；

（2） 若 $p=\dfrac12$，且 $\{a_{2n-1}\}$ 是递增数列，$\{a_{2n}\}$ 是递减数列，求数列 $\{a_n\}$ 的通项公式。

\Needspace{0.4\textheight}
\end{question}
\begin{question}
\gkti{21}如图，$O$ 为坐标原点，椭圆 $C_1:\dfrac{x^2}{a^2}+\dfrac{y^2}{b^2}=1$（$a>b>0$）的左、右焦点分别为 $F_1$，$F_2$，离心率为 $e_1$；双曲线 $C_2:\dfrac{x^2}{a^2}-\dfrac{y^2}{b^2}=1$ 的左、右焦点分别为 $F_3$，$F_4$，离心率为 $e_2$。已知 $e_1e_2=\dfrac{\sqrt3}2$，且 $|F_2F_4|=\sqrt3-1$。
\nopagebreak[4]
\begin{center}% ============================================================
% f2014_hunan_science_q21_conics.tex
% 2014 湖南卷（理科）第 21 题：椭圆 C1 与双曲线 C2
% 来源 PDF：普通高考/2014/2014湖南理.pdf
% 原页：第 2 页（A3 横排）中栏，印面 886，第 21 题题干下方
% 结构：C1: x^2/2+y^2=1；C2: x^2/2-y^2=1；
%       F1(-1,0)、F2(1,0)、F3(-√3,0)、F4(√3,0)；
%       弦 AB 过 F1，M 为 AB 中点，直线 OM 与 C2 交于 P、Q，
%       图中画出四边形 APBQ 的四边。取 AB 的斜率倒数为 m=1/2 作图。
%       比例 1.7 cm/单位。
% ============================================================
\begin{tikzpicture}[>=Stealth, line width=0.6pt, font=\small,
  x=1.7cm, y=1.7cm, line cap=round, line join=round]
  % 坐标轴
  \draw[->] (-2.30,0) -- (2.30,0) node[right] {$x$};
  \draw[->] (0,-1.62) -- (0,1.62) node[above] {$y$};
  % 椭圆 C1
  \draw[domain=0:360, samples=240, smooth, variable=\t]
    plot ({1.41421*cos(\t)}, {sin(\t)});
  % 双曲线 C2 的左右两支
  \draw[domain=-1.50:1.50, samples=120, smooth, variable=\t]
    plot ({sqrt(2*(1+\t*\t))}, {\t});
  \draw[domain=-1.50:1.50, samples=120, smooth, variable=\t]
    plot ({-sqrt(2*(1+\t*\t))}, {\t});
  % 弦 AB（过 F1）与直线 OM（即 PQ）
  \draw (-0.5375,0.9250) -- (-1.2400,-0.4800);
  \draw (-1.5120,0.3780) -- (1.5120,-0.3780);
  % 四边形 APBQ 的四边
  \draw (-0.5375,0.9250) -- (-1.5120,0.3780);
  \draw (-1.5120,0.3780) -- (-1.2400,-0.4800);
  \draw (-1.2400,-0.4800) -- (1.5120,-0.3780);
  \draw (1.5120,-0.3780) -- (-0.5375,0.9250);
  % 焦点
  \fill (-1,0) circle (1.2pt) node[below right, inner sep=2pt] {$F_1$};
  \fill (1,0) circle (1.2pt) node[above, inner sep=2pt] {$F_2$};
  \fill (-1.73205,0) circle (1.2pt) node[below left, inner sep=2pt] {$F_3$};
  \fill (1.73205,0) circle (1.2pt) node[above, inner sep=2pt] {$F_4$};
  % 原点、中点与交点
  \node[below left, inner sep=2pt] at (-0.03,-0.03) {$O$};
  \fill (-0.8889,0.2222) circle (1.0pt);
  \node[above right, fill=white, inner sep=1pt] at (-0.86,0.28) {$M$};
  \node[above, fill=white, inner sep=1pt] at (-1.5120,0.44) {$P$};
  \node[below right, fill=white, inner sep=1pt] at (1.50,-0.42) {$Q$};
  \node[above left, fill=white, inner sep=1pt] at (-0.55,0.93) {$A$};
  \node[below left, fill=white, inner sep=1pt] at (-1.26,-0.50) {$B$};
\end{tikzpicture}
\end{center}
\nopagebreak[4]
（1） 求 $C_1$，$C_2$ 的方程；

（2） 过 $F_1$ 作 $C_1$ 的不垂直于 $y$ 轴的弦 $AB$，$M$ 为弦 $AB$ 的中点。当直线 $OM$ 与 $C_2$ 交于 $P$，$Q$ 两点时，求四边形 $APBQ$ 面积的最小值。

\end{question}
\begin{question}
\gkti{22}已知常数 $a>0$，函数 $f(x)=\ln(1+ax)-\dfrac{2x}{x+2}$。

（1） 讨论 $f(x)$ 在区间 $(0,+\infty)$ 上的单调性；

（2） 若 $f(x)$ 存在两个极值点 $x_1$，$x_2$，且 $f(x_1)+f(x_2)>0$，求 $a$ 的取值范围。

\end{question}